% ===========================================================================
%  Jàng Afrig — Lab 0 : Mise en route
%  Installation de l'environnement de développement
%  Licence 1 DUT1-Informatique — ESP/UCAD — 2025–2026
%  Dr. El Hadji Bassirou TOURÉ
%  Compilé avec : xelatex (2–3 passes)
% ===========================================================================
\documentclass[11pt,a4paper]{article}

\usepackage{fontspec}
\setmainfont{Latin Modern Roman}
\setsansfont{Latin Modern Sans}
\setmonofont{Latin Modern Mono}
\usepackage{polyglossia}
\setdefaultlanguage{french}

\usepackage[margin=2.5cm]{geometry}
\usepackage{amsmath,amssymb}
\usepackage{graphicx}
\graphicspath{{figures/}}
\usepackage{xcolor}
\usepackage{booktabs}
\usepackage{hyperref}
\usepackage{tcolorbox}
\tcbuselibrary{theorems,breakable,skins}
\usepackage{enumitem}
\usepackage{fancyhdr}
\usepackage{array}
\usepackage{pifont}
\usepackage{longtable}
\usepackage{multirow}

% --- TikZ pour les schémas ---
\usepackage{tikz}
\usetikzlibrary{shapes.geometric,arrows.meta,positioning,calc,backgrounds,fit,shadows.blur}

\newcommand{\cmark}{\ding{51}}
\newcommand{\xmark}{\ding{55}}

\hypersetup{
    colorlinks=true,
    linkcolor=blue!60!black,
    urlcolor=blue!60!black,
    citecolor=blue!60!black
}

\newtcolorbox{objectifbox}[1][]{colback=blue!5,colframe=blue!50,fonttitle=\bfseries,title=#1,breakable}
\newtcolorbox{infobox}[1][]{colback=green!5,colframe=green!50!black,fonttitle=\bfseries,title=#1,breakable}
\newtcolorbox{attentionbox}[1][]{colback=orange!5,colframe=orange!70,fonttitle=\bfseries,title=#1,breakable}
\newtcolorbox{retenirbox}[1][]{colback=yellow!10,colframe=yellow!50!black,fonttitle=\bfseries,title=#1,breakable}
\newtcolorbox{prereqbox}[1][]{colback=cyan!5,colframe=cyan!50!black,fonttitle=\bfseries,title=#1,breakable}
\newtcolorbox{projetbox}[1][]{colback=teal!5,colframe=teal!50!black,fonttitle=\bfseries,title=#1,breakable}

% --- Couleurs pour les schémas ---
\definecolor{outilvscode}{RGB}{0,122,204}
\definecolor{outilnav}{RGB}{229,103,23}
\definecolor{outilgit}{RGB}{240,80,50}
\definecolor{outilgithub}{RGB}{36,41,46}
\definecolor{fluxlocal}{RGB}{45,160,110}
\definecolor{fluxdistant}{RGB}{110,70,180}

% --- Code source ---
\usepackage{listings}
\lstset{
  basicstyle=\ttfamily\small,
  backgroundcolor=\color{gray!8},
  frame=single,
  framerule=0.5pt,
  rulecolor=\color{gray!40},
  breaklines=true,
  breakatwhitespace=true,
  showstringspaces=false,
  tabsize=2,
  xleftmargin=4pt,
  xrightmargin=4pt,
  aboveskip=8pt,
  belowskip=8pt,
  literate={é}{{\'e}}1 {è}{{\`e}}1 {ê}{{\^e}}1 {à}{{\`a}}1 {ù}{{\`u}}1
           {ô}{{\^o}}1 {î}{{\^i}}1 {ç}{{\c{c}}}1 {É}{{\'E}}1
           {→}{{$\rightarrow$}}1
}

\pagestyle{fancy}
\fancyhf{}
\fancyhead[L]{\small Jàng Afrig --- Lab 0 : Mise en route}
\fancyhead[R]{\small UCAD --- L1 DUT1-Info}
\fancyfoot[C]{\thepage}
\fancyfoot[L]{\small Dr. El Hadji Bassirou TOURÉ}
\fancyfoot[R]{\small 2025--2026}
\renewcommand{\headrulewidth}{0.4pt}
\renewcommand{\footrulewidth}{0.4pt}

\title{%
\vspace{-1cm}
\textbf{\LARGE Fondamentaux du Web}\\
\textbf{\LARGE HTML, CSS, JavaScript}\\[0.8cm]
\textbf{\huge Lab 0 --- Mise en route}\\[0.3cm]
\textbf{\Large Préparer son environnement de développement}\\[0.5cm]
\large Licence 1 DUT1-Informatique --- 2025--2026}
\author{Dr. El Hadji Bassirou TOURÉ\\
Département Génie Informatique\\
Université Cheikh Anta Diop de Dakar}
\date{}

\begin{document}

\maketitle
\thispagestyle{empty}

\vspace{0.5cm}

% ===========================================================================
% BOÎTE PRÉREQUIS
% ===========================================================================
\begin{prereqbox}[Informations pratiques]
\begin{center}
\begin{tabular}{ll}
\toprule
\textbf{Durée estimée} & 1h30 à 2h (à faire chez soi avant la S1) \\
\textbf{Où} & Chez vous, sur votre ordinateur personnel \\
\textbf{Prérequis} & Bases d'utilisation d'un terminal + connaissance minimale HTML \\
\textbf{Matériel} & Un ordinateur (Windows, macOS ou Linux) + connexion Internet \\
\textbf{Validation} & Un fichier \texttt{index.html} poussé sur votre repo GitHub \\
\textbf{Rapport} & 1 page : ce que j'ai fait + difficultés + lien repo \\
\bottomrule
\end{tabular}
\end{center}
\end{prereqbox}

\begin{objectifbox}[Objectifs du Lab]
À la fin de ce Lab, vous aurez :
\begin{enumerate}[nosep]
    \item Compris comment les \textbf{4 outils} du développeur web (VS Code, navigateur, Git, GitHub) travaillent ensemble
    \item Installé \textbf{VS Code} et l'extension \textbf{Live Server}
    \item Installé et configuré \textbf{Git} avec votre nom et votre email
    \item Créé un compte \textbf{GitHub} et votre dépôt \texttt{jang-afrig}
    \item Cloné ce dépôt sur votre ordinateur
    \item Écrit une première page HTML «\,Hello\,» et l'avoir affichée avec Live Server
    \item Fait votre \textbf{premier cycle Git} : \texttt{add} $\to$ \texttt{commit} $\to$ \texttt{push}
    \item Vérifié votre travail sur GitHub
    \item Rédigé votre premier rapport de Lab
\end{enumerate}
\end{objectifbox}

\tableofcontents

\newpage

% ===========================================================================
\section{Le paysage du développeur web}
% ===========================================================================

Avant de cliquer quoi que ce soit, prenons 5 minutes pour comprendre
\textbf{avec quels outils vous allez travailler} tout le semestre. Un
développeur web travaille avec \textbf{quatre outils} qui se complètent,
comme un artisan dans son atelier :

\begin{itemize}[nosep]
    \item \textbf{Un éditeur de code} --- VS Code, votre \emph{table de travail}
    \item \textbf{Un navigateur} --- Chrome ou Firefox, votre \emph{vitrine}
    \item \textbf{Un outil de versioning} --- Git, votre \emph{carnet de bord}
    \item \textbf{Un hébergeur} --- GitHub, votre \emph{entrepôt partagé}
\end{itemize}

Voici comment ils travaillent ensemble :

\begin{figure}[!ht]
\centering
\begin{tikzpicture}[
    node distance=2cm,
    outil/.style={rounded corners=6pt, draw, thick, minimum width=3.2cm, minimum height=1.8cm, align=center, font=\small\bfseries, drop shadow={shadow xshift=1pt, shadow yshift=-1pt, opacity=0.3}},
    fleche/.style={-{Stealth[length=3mm]}, thick, shorten >=2pt, shorten <=2pt},
    libelle/.style={font=\scriptsize\itshape, midway, fill=white, inner sep=2pt}
]
    % Les 4 outils
    \node[outil, fill=outilvscode!15, draw=outilvscode] (vscode) {VS Code\\ \normalfont\small Écrire le code};
    \node[outil, fill=outilnav!15, draw=outilnav, right=3cm of vscode] (navigateur) {Navigateur\\ \normalfont\small Afficher la page};
    \node[outil, fill=outilgit!15, draw=outilgit, below=2cm of vscode] (git) {Git\\ \normalfont\small Sauvegarder les versions};
    \node[outil, fill=outilgithub!20, draw=outilgithub, below=2cm of navigateur] (github) {GitHub\\ \normalfont\small Stocker en ligne};

    % Flèches
    \draw[fleche] (vscode) -- node[libelle, above] {j'écris} (navigateur);
    \draw[fleche] (navigateur) -- node[libelle, below] {je vois, je teste} (vscode);

    \draw[fleche] (vscode) -- node[libelle, right] {j'enregistre} (git);
    \draw[fleche] (git) -- node[libelle, below] {je pousse} (github);
\end{tikzpicture}
\caption{Les quatre outils du développeur web et leurs interactions.}
\label{fig:outils}
\end{figure}

\begin{retenirbox}[À retenir]
\begin{itemize}[nosep]
    \item J'\textbf{écris} mon code dans VS Code et je le \textbf{teste} dans le navigateur
    \item J'\textbf{enregistre} mes versions avec Git sur mon ordinateur
    \item Je \textbf{pousse} mes versions sur GitHub pour les mettre en ligne et les sauvegarder
\end{itemize}
\end{retenirbox}

\newpage

% ===========================================================================
\section{Ce qu'on va installer et pourquoi}
% ===========================================================================

\begin{center}
\begin{tabular}{@{}lp{8cm}@{}}
\toprule
\textbf{Outil} & \textbf{À quoi il sert} \\
\midrule
\textbf{VS Code} & L'éditeur de code. C'est là que vous écrirez HTML, CSS, JS. Intelligent : il colore les mots-clés, détecte les erreurs, complète automatiquement. \\
\addlinespace
\textbf{Live Server} & Extension de VS Code. Lance un petit serveur local et \textbf{recharge automatiquement} la page dans le navigateur dès que vous enregistrez le code. Gain de temps énorme. \\
\addlinespace
\textbf{Git} & Outil de \emph{versioning}. Enregistre l'historique de vos modifications. Vous pouvez revenir en arrière à n'importe quel moment. C'est un gros bouton «\,annuler\,» sans limite. \\
\addlinespace
\textbf{GitHub} & Site web qui héberge vos dépôts Git en ligne. Comme Google Drive, mais pour les développeurs. Vous y déposerez tout votre projet Jàng Afrig. \\
\bottomrule
\end{tabular}
\end{center}

\begin{attentionbox}[Règle d'or]
Ne passez jamais à l'étape suivante tant que l'étape actuelle n'a pas
produit le résultat attendu. Si ça bloque, relisez calmement, cherchez la
coquille. \textbf{Chercher par soi-même fait partie du métier.}
\end{attentionbox}

% ===========================================================================
\section{Installer VS Code}
% ===========================================================================

\subsection{Télécharger et installer}

\begin{enumerate}[nosep]
    \item Allez sur \url{https://code.visualstudio.com/}
    \item Cliquez sur le bouton \textbf{Download} correspondant à votre système (Windows, macOS, Linux)
    \item Lancez le fichier téléchargé et suivez l'assistant
    \item Gardez les options par défaut ; sur Windows, \textbf{cochez «\,Ajouter à PATH\,»}
\end{enumerate}

\begin{infobox}[Astuce VS Code]
VS Code permet de faire énormément de choses via la palette de commandes :
\texttt{Ctrl+Shift+P} (Windows/Linux) ou \texttt{Cmd+Shift+P} (macOS). On
s'en servira plusieurs fois dans le semestre.
\end{infobox}

\subsection{Installer l'extension Live Server}

\begin{enumerate}[nosep]
    \item Ouvrez VS Code
    \item Cliquez sur l'icône \textbf{Extensions} dans la barre de gauche (4 carrés)
    \item Dans la barre de recherche, tapez : \texttt{Live Server}
    \item Choisissez celui de \textbf{Ritwick Dey} (le plus téléchargé)
    \item Cliquez sur \textbf{Install}
\end{enumerate}

\begin{retenirbox}[Pourquoi Live Server ?]
Sans Live Server, à chaque modification de code, vous devez :
enregistrer $\to$ aller dans le navigateur $\to$ rafraîchir la page.
Avec Live Server : vous enregistrez, la page se rafraîchit \textbf{toute
seule}. Sur 30h de cours, ce petit détail vous économise des heures.
\end{retenirbox}

\newpage

% ===========================================================================
\section{Installer et configurer Git}
% ===========================================================================

\subsection{Installer Git}

\subsubsection*{Windows}
\begin{enumerate}[nosep]
    \item Allez sur \url{https://git-scm.com/download/win}
    \item Le téléchargement démarre automatiquement
    \item Lancez le fichier, cliquez \textbf{Next} partout sauf :
    \item Quand on vous demande l'éditeur par défaut, choisissez \textbf{«\,Use Visual Studio Code\,»}
\end{enumerate}

\subsubsection*{macOS}
\begin{lstlisting}
git --version
\end{lstlisting}
Si Git n'est pas installé, macOS vous proposera de l'installer (acceptez).

\subsubsection*{Linux (Ubuntu/Debian)}
\begin{lstlisting}
sudo apt update
sudo apt install git
\end{lstlisting}

\subsection{Vérifier l'installation}

Ouvrez votre terminal (Git Bash sous Windows, Terminal sous macOS/Linux)
et tapez :

\begin{lstlisting}
git --version
\end{lstlisting}

Résultat attendu :

\begin{lstlisting}
git version 2.47.1
\end{lstlisting}

\cmark{} Si vous voyez un numéro de version, Git est installé.

\subsection{Configurer Git avec votre identité}

Git a besoin de savoir \textbf{qui vous êtes}. Chaque enregistrement portera
votre nom et votre email. Tapez ces deux commandes (avec \textbf{votre}
vrai nom et \textbf{votre} email étudiant) :

\begin{lstlisting}
git config --global user.name "Fatou Diop"
git config --global user.email "fatou.diop@esp.sn"
\end{lstlisting}

\begin{attentionbox}[Utilisez votre vrai nom et email]
Remplacez par votre vrai nom (pas un pseudo) et votre email étudiant. Ces
informations apparaîtront dans l'historique Git et servent à l'évaluation.
\end{attentionbox}

Vérifiez :
\begin{lstlisting}
git config --global user.name
git config --global user.email
\end{lstlisting}

Vous devez voir votre nom et email s'afficher. \cmark

\newpage

% ===========================================================================
\section{Créer votre compte GitHub}
% ===========================================================================

\subsection{S'inscrire}

\begin{enumerate}[nosep]
    \item Allez sur \url{https://github.com/}
    \item Cliquez sur \textbf{Sign up}
    \item Utilisez le même email que pour Git
    \item Choisissez un mot de passe solide
    \item Choisissez un \textbf{nom d'utilisateur} (\emph{username})
    \item Complétez l'inscription et vérifiez votre email
\end{enumerate}

\begin{attentionbox}[Choisissez un username professionnel]
Votre \emph{username} GitHub apparaîtra sur votre CV pendant toute votre
carrière. Évitez :
\begin{itemize}[nosep]
    \item \texttt{fatoudupays221} \xmark{} (trop informel)
    \item \texttt{xxx\_killer\_xxx} \xmark{} (à ne pas montrer à un recruteur)
\end{itemize}
Préférez :
\begin{itemize}[nosep]
    \item \texttt{fatou-diop} \cmark
    \item \texttt{fdiop} \cmark
    \item \texttt{fatou.diop} \cmark
\end{itemize}
\end{attentionbox}

\subsection{Créer votre dépôt jang-afrig}

\begin{enumerate}[nosep]
    \item Connectez-vous sur GitHub
    \item Cliquez sur le \textbf{+} en haut à droite $\to$ \textbf{New repository}
    \item \textbf{Repository name} : \texttt{jang-afrig}
    \item \textbf{Description} : \texttt{Quiz géographique sur les 54 pays africains}
    \item Laissez \textbf{Public}
    \item Cochez \textbf{«\,Add a README file\,»}
    \item Cliquez \textbf{Create repository}
\end{enumerate}

\subsection{Récupérer l'URL du dépôt}

Sur la page de votre dépôt, cliquez sur le bouton vert \textbf{Code}, onglet
\textbf{HTTPS}, copiez l'URL qui finit par \texttt{.git}. Elle ressemble à :

\begin{lstlisting}
https://github.com/votre-username/jang-afrig.git
\end{lstlisting}

Gardez-la de côté, on en a besoin à l'étape suivante.

\newpage

% ===========================================================================
\section{Cloner le dépôt sur votre ordinateur}
% ===========================================================================

\textbf{Cloner} un dépôt, c'est faire une copie locale sur votre ordinateur.
Vous pourrez éditer les fichiers chez vous, puis renvoyer vos modifications
sur GitHub.

\subsection{Créer un dossier pour vos projets}

Dans le terminal, allez dans votre dossier personnel et créez un dossier
dédié à ce cours :

\begin{lstlisting}
cd ~
mkdir cours-web
cd cours-web
\end{lstlisting}

\subsection{Cloner le dépôt}

Utilisez la commande \texttt{git clone} avec votre URL (en remplaçant
\texttt{votre-username}) :

\begin{lstlisting}
git clone https://github.com/votre-username/jang-afrig.git
\end{lstlisting}

Résultat attendu :

\begin{lstlisting}
Cloning into 'jang-afrig'...
remote: Enumerating objects: 3, done.
Receiving objects: 100% (3/3), done.
\end{lstlisting}

Un dossier \texttt{jang-afrig} a été créé. Entrez dedans :

\begin{lstlisting}
cd jang-afrig
ls
\end{lstlisting}

Vous devez voir au moins un fichier \texttt{README.md}.

\newpage

% ===========================================================================
\section{Créer votre première page HTML}
% ===========================================================================

\subsection{Ouvrir le dossier dans VS Code}

Dans le terminal, toujours dans le dossier \texttt{jang-afrig}, tapez :

\begin{lstlisting}
code .
\end{lstlisting}

\textbf{Attention au point.} Cette commande ouvre le dossier courant dans
VS Code.

\begin{infobox}[Si \texttt{code .} ne marche pas sur macOS]
Dans VS Code, \texttt{Cmd+Shift+P} $\to$ tapez «\,shell command\,» $\to$
«\,Install 'code' command in PATH\,». Sinon, ouvrez VS Code et faites
\textbf{File $\to$ Open Folder}.
\end{infobox}

\subsection{Créer index.html}

\begin{enumerate}[nosep]
    \item Dans l'Explorateur de VS Code (barre de gauche), cliquez sur l'icône «\,Nouveau fichier\,» à côté de \textbf{NJURELGEO}
    \item Tapez \texttt{index.html}
    \item Copiez-collez ce code (remplacez «\,Fatou\,» par \textbf{votre prénom}) :
\end{enumerate}

\begin{lstlisting}
<!DOCTYPE html>
<html lang="fr">
<head>
    <meta charset="UTF-8">
    <title>Hello Fatou</title>
</head>
<body>
    <h1>Hello Fatou !</h1>
    <p>Bienvenue dans mon premier projet web : Jàng Afrig.</p>
</body>
</html>
\end{lstlisting}

Enregistrez avec \texttt{Ctrl+S} ou \texttt{Cmd+S}.

\subsection{Afficher avec Live Server}

\begin{enumerate}[nosep]
    \item Clic droit sur \texttt{index.html} dans l'Explorateur
    \item \textbf{Open with Live Server}
\end{enumerate}

Votre navigateur s'ouvre sur la page. \cmark

\begin{infobox}[Testez le rechargement automatique]
Dans VS Code, changez «\,Hello Fatou\,» en «\,Salaam Aleykum Fatou\,» et
enregistrez. La page se met à jour \textbf{toute seule} dans le navigateur.
C'est ça la magie de Live Server.
\end{infobox}

\newpage

% ===========================================================================
\section{Premier cycle Git : du local vers GitHub}
% ===========================================================================

Vous avez créé un nouveau fichier sur votre ordinateur. Il faut maintenant
l'\textbf{enregistrer dans l'historique Git}, puis l'\textbf{envoyer sur
GitHub}. C'est ce qu'on appelle un \emph{cycle Git}.

Le schéma ci-dessous montre le flux à comprendre :

\begin{figure}[!ht]
\centering
\begin{tikzpicture}[
    node distance=1.3cm,
    zone/.style={rounded corners=6pt, draw, thick, minimum width=4.8cm, minimum height=1.1cm, align=center, font=\small\bfseries},
    fleche/.style={-{Stealth[length=3mm]}, thick, shorten >=3pt, shorten <=3pt, line width=1pt},
    cmd/.style={font=\ttfamily\small, fill=yellow!15, rounded corners=3pt, inner sep=4pt, draw=yellow!60!black, thick}
]
    % Les 4 zones en vertical
    \node[zone, fill=fluxlocal!10, draw=fluxlocal] (wd) {Fichiers sur mon ordinateur};
    \node[zone, fill=fluxlocal!15, draw=fluxlocal, below=of wd] (stage) {Zone de préparation};
    \node[zone, fill=fluxlocal!20, draw=fluxlocal, below=of stage] (local) {Historique local};
    \node[zone, fill=fluxdistant!15, draw=fluxdistant, below=of local] (remote) {GitHub (en ligne)};

    % Commandes entre les zones
    \draw[fleche] (wd) -- (stage);
    \node[cmd, right=0.1cm of $(wd)!0.5!(stage)$] {git add};

    \draw[fleche] (stage) -- (local);
    \node[cmd, right=0.1cm of $(stage)!0.5!(local)$] {git commit};

    \draw[fleche] (local) -- (remote);
    \node[cmd, right=0.1cm of $(local)!0.5!(remote)$] {git push};

    % Légende des étapes (à gauche)
    \node[left=0.3cm of wd, font=\scriptsize\itshape, text=gray] {j'édite};
    \node[left=0.3cm of stage, font=\scriptsize\itshape, text=gray] {je prépare};
    \node[left=0.3cm of local, font=\scriptsize\itshape, text=gray] {j'enregistre};
    \node[left=0.3cm of remote, font=\scriptsize\itshape, text=gray] {je partage};

    % Frame local vs distant
    \begin{scope}[on background layer]
        \node[fit=(wd)(stage)(local), rounded corners=10pt, draw=fluxlocal!60, dashed, thick, inner sep=8pt, label={[font=\small\bfseries, text=fluxlocal]right:Sur mon ordinateur}] {};
        \node[fit=(remote), rounded corners=10pt, draw=fluxdistant!60, dashed, thick, inner sep=8pt, label={[font=\small\bfseries, text=fluxdistant]right:En ligne}] {};
    \end{scope}
\end{tikzpicture}
\caption{Le flux d'un cycle Git : du fichier local jusqu'à GitHub.}
\label{fig:flux-git}
\end{figure}

\begin{retenirbox}[L'analogie du colis]
\begin{itemize}[nosep]
    \item Vous avez un cadeau chez vous (votre \textbf{fichier}).
    \item Vous le mettez dans un carton (\texttt{git add} $\to$ zone de préparation).
    \item Vous fermez le carton et lui mettez une étiquette avec un message (\texttt{git commit}).
    \item Vous donnez le carton à La Poste pour qu'il parte à l'étranger (\texttt{git push} $\to$ GitHub).
\end{itemize}
Chaque carton (\emph{commit}) reste dans l'historique. Vous pouvez toujours
le retrouver, même 6 mois plus tard.
\end{retenirbox}

\newpage

\subsection{Ouvrir le terminal dans VS Code}

VS Code a un terminal intégré. Pour l'ouvrir :
\begin{itemize}[nosep]
    \item Menu \textbf{Terminal $\to$ New Terminal}
    \item Raccourci : \texttt{Ctrl+\textasciigrave} / \texttt{Cmd+\textasciigrave}
\end{itemize}

Le terminal s'ouvre en bas, déjà placé dans \texttt{jang-afrig}.

\subsection{Voir l'état avec git status}

\begin{lstlisting}
git status
\end{lstlisting}

Résultat :
\begin{lstlisting}
On branch main
Your branch is up to date with 'origin/main'.

Untracked files:
  (use "git add <file>..." to include in what will be committed)
        index.html

nothing added to commit but untracked files present
\end{lstlisting}

Git dit : «\,Je vois un nouveau fichier, mais je ne le surveille pas encore.\,»

\subsection{Préparer avec git add}

\begin{lstlisting}
git add index.html
\end{lstlisting}

Pas de message : Git est silencieux quand tout va bien. Vérifiez :

\begin{lstlisting}
git status
\end{lstlisting}

\begin{lstlisting}
Changes to be committed:
        new file:   index.html
\end{lstlisting}

Le fichier est \textbf{préparé}.

\subsection{Enregistrer avec git commit}

\begin{lstlisting}
git commit -m "v0 : premier setup avec hello world"
\end{lstlisting}

Le texte entre guillemets est un \textbf{message de commit}. Décrivez
brièvement ce que vous venez de faire.

\begin{lstlisting}
[main a1b2c3d] v0 : premier setup avec hello world
 1 file changed, 10 insertions(+)
\end{lstlisting}

\cmark{} Commit enregistré en local.

\subsection{Envoyer sur GitHub avec git push}

\begin{lstlisting}
git push
\end{lstlisting}

La première fois, Git vous demandera votre \textbf{username} GitHub et un
\textbf{token}.

\begin{attentionbox}[GitHub : utilisez un Personal Access Token (PAT)]
GitHub n'accepte plus les mots de passe depuis 2021. À la place, vous avez
besoin d'un \textbf{Personal Access Token} (PAT) :

\begin{enumerate}[nosep]
    \item GitHub $\to$ Settings $\to$ Developer settings $\to$ Personal access tokens $\to$ Tokens (classic)
    \item \textbf{Generate new token (classic)}
    \item Nom : «\,cours-web-2026\,»
    \item Cochez la case \textbf{repo}
    \item \textbf{Generate} $\to$ \textbf{COPIEZ LE TOKEN TOUT DE SUITE} --- vous ne le reverrez plus
    \item Utilisez ce token à la place du mot de passe
\end{enumerate}

Notez-le dans un endroit sûr (gestionnaire de mots de passe).
\end{attentionbox}

Résultat attendu :
\begin{lstlisting}
Enumerating objects: 4, done.
Writing objects: 100% (3/3), 394 bytes
To https://github.com/votre-username/jang-afrig.git
   1a2b3c4..a1b2c3d  main -> main
\end{lstlisting}

\cmark{} Votre code est sur GitHub.

\subsection{Vérifier sur GitHub}

Allez sur \url{https://github.com/votre-username/jang-afrig}. Vous devez
voir \texttt{index.html} dans la liste. Cliquez dessus pour voir son contenu.

\newpage

% ===========================================================================
\section{Aide-mémoire}
% ===========================================================================

\subsection{Les commandes apprises dans ce Lab}

\begin{center}
\begin{tabular}{lp{8cm}}
\toprule
\textbf{Commande} & \textbf{Ce qu'elle fait} \\
\midrule
\texttt{git --version} & Vérifier que Git est installé \\
\texttt{git config --global user.name "..."} & Définir votre nom Git \\
\texttt{git config --global user.email "..."} & Définir votre email Git \\
\texttt{git clone <url>} & Télécharger un dépôt depuis GitHub \\
\texttt{git status} & Voir l'état des fichiers \\
\texttt{git add <fichier>} & Préparer un fichier \\
\texttt{git commit -m "message"} & Enregistrer les fichiers préparés \\
\texttt{git push} & Envoyer sur GitHub \\
\texttt{code .} & Ouvrir le dossier actuel dans VS Code \\
\bottomrule
\end{tabular}
\end{center}

\subsection{Le cycle Git : la routine du semestre}

\begin{center}
\begin{tabular}{rl}
1. & Modifier des fichiers dans VS Code \\
2. & \texttt{git status} \quad (voir ce qui a changé) \\
3. & \texttt{git add .} \quad (préparer tous les changements) \\
4. & \texttt{git commit -m "message"} \quad (enregistrer en local) \\
5. & \texttt{git push} \quad (envoyer sur GitHub) \\
\end{tabular}
\end{center}

\begin{retenirbox}[À retenir absolument]
Ces 5 étapes seront votre routine à la fin de chaque Lab pendant tout le
semestre. Vous allez les faire 10 fois. Au bout de 3 ou 4, vous les ferez
sans réfléchir.
\end{retenirbox}

\subsection{Si quelque chose ne marche pas}

\begin{infobox}[Méthode de déblocage]
\begin{enumerate}[nosep]
    \item Relisez votre commande --- une coquille de lettre bloque tout
    \item Copiez le message d'erreur exact dans un moteur de recherche
    \item Demandez à un camarade qui a réussi
    \item Notez le problème et montrez-le au prof au début de la S1
\end{enumerate}
\end{infobox}

\newpage

% ===========================================================================
\section{Rapport à produire}
% ===========================================================================

\begin{attentionbox}[Rapport du Lab 0 --- obligatoire]
À la fin de ce Lab, vous produisez un \textbf{rapport court} de
\textbf{1 à 2 pages maximum}. Ce sera la routine pour chaque Lab du semestre.

\medskip

\textbf{Format du fichier :}
\begin{itemize}[nosep]
    \item PDF, Markdown (\texttt{.md}) ou Word
    \item Nom : \texttt{rapport-lab-0-NOM-Prenom.pdf}\\
          (ex. \texttt{rapport-lab-0-DIOP-Fatou.pdf})
    \item Déposer dans un sous-dossier \texttt{rapports/} de votre repo GitHub
\end{itemize}

\medskip

\textbf{Contenu obligatoire (4 sections) :}

\begin{enumerate}[nosep]
    \item \textbf{Ce que j'ai fait} (½ page) : outils installés, configuration Git, création du repo, premier cycle Git.
    \item \textbf{Difficultés rencontrées} (¼ page) : ce qui m'a bloqué, messages d'erreur rencontrés.
    \item \textbf{Comment j'ai résolu ces difficultés} (¼ page) : recherche, aide, essais.
    \item \textbf{Lien vers mon repo GitHub} : \\
          ex. \url{https://github.com/fatou-diop/jang-afrig}
\end{enumerate}

\medskip

\textbf{Ton attendu :} honnête et direct. Si quelque chose n'a pas marché
du premier coup, dites-le. L'erreur fait partie de l'apprentissage et
\textbf{personne ne vous jugera dessus}. Ce qui compte, c'est que vous
ayez réfléchi à ce que vous avez fait.
\end{attentionbox}

% ===========================================================================
\section{Résumé et conclusion}
% ===========================================================================

Dans ce Lab, vous avez :

\begin{enumerate}[nosep]
    \item Compris le rôle des 4 outils et leurs interactions \cmark
    \item Installé VS Code et Live Server \cmark
    \item Installé et configuré Git avec votre identité \cmark
    \item Créé un compte GitHub et votre premier dépôt \cmark
    \item Cloné ce dépôt sur votre ordinateur \cmark
    \item Écrit votre première page HTML avec Live Server \cmark
    \item Fait votre premier cycle Git : \texttt{add} $\to$ \texttt{commit} $\to$ \texttt{push} \cmark
    \item Rédigé votre premier rapport de Lab \cmark
\end{enumerate}

\textbf{Votre environnement de développement est prêt.} Vous êtes paré(e)
pour la S1.

\subsection{Et après ?}

Au CM1 (S1), nous allons :

\begin{itemize}[nosep]
    \item Faire le tour de ce qu'est vraiment une page web
    \item Découvrir Jàng Afrig dans le détail
    \item Approfondir les balises HTML
    \item Commencer à construire la page d'accueil
\end{itemize}

\begin{projetbox}[Votre vision à 15 séances]
Votre dépôt \texttt{jang-afrig} contient aujourd'hui une simple page
«\,Hello\,». Dans 15 séances, il contiendra une vraie application de quiz
géographique, en ligne, accessible depuis votre téléphone, utilisant les
données de 54 pays africains récupérées sur une API.

Vous pourrez la montrer à votre famille via un simple lien.
\end{projetbox}

\vspace{1cm}

\begin{center}
\fbox{
\begin{minipage}{0.85\textwidth}
\centering
\vspace{0.5cm}
\textbf{\Large Lab 0 --- Terminé \cmark}\\[0.3cm]
\normalsize
Votre atelier est prêt. Bienvenue dans le monde du développement web.\\[0.2cm]
\textit{Rendez-vous au CM1 !}\\[0.3cm]
\vspace{0.3cm}
\end{minipage}
}
\end{center}

\end{document}
